% AG.32--AG.33: full direct proof; independently approved by root and criteria.
\begin{claim}{AG.32}{Proved; an effective conductor bound in the scalar supersingular case}
Let $C/\mathbf F_{2^{2r}}$ be a genus-two curve with two-rank zero,
with rational Weierstrass point $P$, and suppose Jacobian Frobenius
is $[\varepsilon2^r]$, where $\varepsilon\in\{1,-1\}$.
If a nontrivial connected cyclic \'etale cover of odd degree $n$
has primitive part of two-rank zero, put
\[
 f=\operatorname{ord}_n(2),\qquad d=\gcd(f,r),\qquad
 M=\frac{\varphi(n)}d.
\]
Then
\[
             M(M+1)\le 2(n-1),\qquad n<10r^3.       \tag{AG.32.1}
\]
Consequently the geometric primitive-cover question for this
fixed scalar-Frobenius curve has only finitely many possible
conductors, with the displayed explicit bound.
\end{claim}
\begin{proof}
The Raynaud theta divisor is the symmetric Abel curve
$\Theta=C-P$, and contains zero. In characteristic two its
nontrivial character cohomology has dimension $g-1=1$.
The orbit criterion AG.3 therefore says that primitive
rank zero requires at least one theta point in every
multiplication-by-two orbit on the primitive characters
of the cyclic subgroup of order $n$.
There are $\varphi(n)/f$ such orbits, each of size $f$.

Frobenius preserves $\Theta$, and symmetry removes the sign
$\varepsilon$. Thus multiplication by $2^r$ preserves its
geometric points. A theta point in one of the required
orbits brings with it its entire multiplication-by-$2^r$
orbit, of size $f/d$. The primitive theta intersection
therefore has at least $M=\varphi(n)/d$ points.
Together with zero there are at least $M+1$ points of
$\Theta$ in the cyclic subgroup. For any nonzero difference
there are at most two ordered pairs of distinct points of
a genus-two Abel curve, since two distinct translates of
its principal theta divisor have intersection number two.
The exact difference count used in AG.27 now gives
$M(M+1)\le2(n-1)$.

For completeness, the elementary uniform estimate
\[
                  \varphi(n)\ge\frac23 n^{5/6}
                       \quad(n\text{ odd})             \tag{AG.32.2}
\]
follows prime-power by prime-power. For primes $\ell\ge5$,
$(\ell-1)^6\ge\ell^5$: it holds at five, and
$(x-1)^6/x^5$ increases for $x>1$. Hence
$\varphi(\ell^a)^6/\ell^{5a}\ge1$ for every $a\ge1$.
The only remaining prime, three, contributes a factor at
least $64/243$, which is larger than $(2/3)^6$.
Multiplying proves (AG.32.2). Since $d\le r$, the first
inequality of (AG.32.1) implies
\[
 \frac{4}{9r^2}n^{5/3}\le M^2<2n,
 \qquad
 n<\left(\frac92\right)^{3/2}r^3<10r^3.
\]
The strict inequality is preserved because $n>1$.
\end{proof}

\begin{claim}{AG.33}{Refuted: universal nontrivial primitive-cover existence}
There are supersingular genus-two curves over
$\mathbf F_{2^{128}}$ for which every nontrivial connected
cyclic \'etale cover of odd degree has primitive part of
positive two-rank. Thus the version of the existence statement
A6 requiring degree $n>1$ is false, already for genus two
in characteristic two.

A completely specified family containing these examples is
\[
 C_t:\qquad y^2+y=
       \left(\frac{u^2+t^2u+t}{t^3+1}\right)^3,
 \qquad
 t\in\mathbf F_{2^{128}}\setminus\mathbf F_4,       \tag{AG.33.1}
\]
where $C_t$ means the smooth proper model. The proof gives
an explicit finite incidence set of excluded parameters
and a positive exact lower bound for its complement.
The program certifies this nonemptiness and the displayed
algebraic identities; it does not claim to have enumerated
one particular parameter in the complement.

All these curves themselves satisfy A1: their Jacobians
have geometric two-rank zero. The refutation concerns
existence of nontrivial primitive-zero covers, and gives no
counterexample to A1 or to the arithmetic criterion A3.
The prompt does not explicitly require $n>1$. With its
convention $P_1=J_C$, the identity cover is a witness for
every curve constructed here. Therefore these examples
do not refute the literal degree-one-permitted version
of A6; that version remains open for curves of positive
$p$-rank.
\end{claim}
\begin{proof}
We first establish the family and its torsion incidence
bound directly. It is the characteristic-two pencil
studied by Moret--Bailly in \emph{Polarisations de degr\'e 4
sur les surfaces ab\'eliennes}, the theorem, parts~1(c) and~3
\citep{AG_MoretBailly1979}; the explicit calculations below
supply the required properties independently.

\emph{1. Explicit genus-two curves and elliptic maps.}
Let $k$ be a finite field containing $\mathbf F_4$, put
$E:y^2+y=x^3$, $A=E\times E$, and take
$t\in k\setminus\mathbf F_4$. Write $d=t^3+1$, so
$t\ne0$ and $d\ne0$. If $X=u^2+t^2u+t$, expansion gives
\[
 X^3=u^6+t^2u^5+(t+t^4)u^4+t^6u^3
                +(t^5+t^2)u^2+t^4u+t^3.
\]
Over the perfect field $k$, successive Artin--Schreier
changes remove the even powers $u^6,u^4,u^2$. The highest
remaining odd power is $u^5$, with nonzero coefficient
$t^2/d^3$. Thus (AG.33.1) has genus two and two-rank zero,
and has exactly one point $P$ above infinity; that point
is $k$-rational and is its unique Weierstrass point.
In particular $\omega_C\simeq\mathcal O_C(2P)$.

Define two morphisms to $E$ by
\[
\begin{aligned}
 x_1&=\frac{u^2+t^2u+t}{d}, & y_1&=y,\\
 x_2&=\frac{tu^2+u+t^2}{d}, & y_2&=y+\frac{u^3+1}{d}.
\end{aligned}                                                   \tag{AG.33.2}
\]
The first elliptic equation is the defining equation of
$C_t$. For the second, put $Z=tu^2+u+t^2$ and $Y=u^3+1$.
The polynomial identity
\[
                 X^3+Z^3=dY^2+d^2Y
\]
proves $y_2^2+y_2=x_2^3$. Properness extends both maps
across $P$, and both take $P$ to the elliptic origin.
The joint map $i_t:C_t\to A$ is birational onto its image,
because
\[
                         u=x_2+t x_1,
                         \qquad y=y_1.                \tag{AG.33.3}
\]

\emph{2. The pencil and its only geometric base point.}
On $A$ consider
\[
 f=\frac{y_2+y_1+1+x_1x_2}{x_2+x_1^2},\qquad
 D_\infty=E\times\{0\}+\Gamma_{F_E},                 \tag{AG.33.4}
\]
where $F_E:E\to E$ is the degree-two Frobenius morphism.
The polar divisor of $f$ is exactly $D_\infty$.
Here is the local check. In the affine part, the denominator
vanishes on the two graphs $\Gamma_{F_E}$ and
$\Gamma_{-F_E}$. On the first graph its numerator is one,
so the pole is simple. On the second the numerator
vanishes and the factor cancels. Along $E\times\{0\}$,
the orders of $y_2$ and $x_2$ are minus three and minus two,
so again the pole is simple with nonzero leading coefficient.
Along $\{0\}\times E$, the numerator has order minus three
and the denominator minus four, so $f$ vanishes rather
than having a pole. These account for every possible
polar divisor.

The two components of $D_\infty$ meet geometrically only
at the origin: the kernel of $F_E$ has only its
identity as a geometric point. At every nonzero point of
the components of $D_\infty$, the preceding nonzero pole
coefficient shows that the section $f$ of
$\mathcal O_A(D_\infty)$ does not vanish. Hence the pencil
generated by $1$ and $f$ has no nonzero geometric base point.
It has the origin as a base point, because (AG.33.2) sends
$P$ to zero for every admissible $t$.

Substitution in (AG.33.4) gives $f\circ i_t=t$.
The divisor $D_t=(f-t)+D_\infty$ is precisely the image
$i_t(C_t)$, with multiplicity one. To check that no
component is being suppressed, write $x=x_1,z=x_2$.
Eliminating $y_2$ from the numerator equation for $f=t$
and the two elliptic equations gives
\[
 (z+x^2)\bigl(z^2+t^2z+t^2x^2+x+t\bigr)=0.
\]
The first factor is the cancelled graph. The second is
exactly (AG.33.1) after $u=z+tx$, as (AG.33.2)--(AG.33.3)
show. It is irreducible because its smooth proper model
has the Artin--Schreier equation with nonzero reduced
$u^5$ coefficient. It has multiplicity one: away from
the denominator, differentiating on $A$ gives
\[
                   df=dx+
                  \frac{z^2+x+f}{z+x^2}\,dz,
\]
which is nonzero along the generic point of this component.
The boundary component $\{0\}\times E$ is a zero fibre
only for $t=0$, excluded here, and the other boundary
component is polar. Therefore $D_t\sim D_\infty$.

\emph{3. The induced isogeny and the theta origin.}
The elliptic components in $D_\infty$ have self-intersection
zero and intersect with length $\deg F_E=2$. Consequently
\[
                         D_t^2=4.                     \tag{AG.33.5}
\]
The map $i_t$, based at $P$, induces a homomorphism
$v_t:J_{C_t}\to A$. Its image is all of $A$: otherwise
the genus-two image curve would lie in a translate of an
elliptic subvariety and would itself have normalization
genus one. Thus $v_t$ is an isogeny.
For the principal Abel theta divisor $\Theta=C_t-P$,
birationality gives $(v_t)_*\Theta=D_t$.
For an isogeny of abelian surfaces, pullback and pushforward
on the rational N\'eron--Severi group satisfy
$v_t^*(v_t)_*=(\deg v_t)\operatorname{id}$, and pullback
multiplies intersection numbers by $\deg v_t$.
Using $\Theta^2=2$ therefore gives
\[
                       D_t^2=2\deg v_t.
\]
Equation (AG.33.5) proves $\deg v_t=2$.
In particular it induces an isomorphism on geometric
$n$-torsion for every odd $n$.

There is no unspecified translation of Raynaud theta in
this argument. Since $C_t$ has two-rank zero,
$J_{C_t}[2](\overline k)=0$. Its theta characteristic is
therefore unique; as $\omega_C=\mathcal O_C(2P)$, the
Raynaud theta characteristic is $\mathcal O_C(P)$.
The corresponding theta divisor is exactly $C_t-P$.
Thus a nonzero theta point killed by an odd $n$ maps to
a nonzero point of $A[n]$ on $D_t$.

\emph{4. Uniform torsion incidence and scalar Frobenius.}
Every nonzero geometric point of $A$ belongs to at most
one member of the pencil, by its single base point.
Since $\#A[n](\overline k)=n^4$ for odd $n$, there are
at most $n^4-1$ parameters over $\overline{\mathbf F}_2$
for which $\Theta_{C_t}$ contains a nonzero $n$-torsion point.
This bounds all geometric torsion points, regardless of
which extension makes them rational.

For $k=\mathbf F_{4^e}$, Frobenius on $E$ is $[(-2)^e]$:
indeed $\#E(\mathbf F_2)=3$, so its Frobenius
polynomial is $X^2+2$, and Cayley--Hamilton for the
endomorphism gives $F_E^2=[-2]$.
The isogeny $v_t$ is defined over $k$, so the same exact
scalar formula holds on $J_{C_t}$. The conductor bound
AG.32 applies with $r=e$. If a primitive-zero cyclic cover
exists, its odd conductor therefore satisfies $3\le n<10e^3$,
and its theta criterion supplies a nonzero theta point
of exact order $n$.

\emph{5. A finite, nonempty counterexample set.}
Take $e=64$, $q=4^{64}=2^{128}$, and $B=10\cdot64^3=2621440$.
Define a finite subset of $\mathbf F_q$ by
\[
 \mathcal B=\mathbf F_4\ \cup\
 \left(\mathbf F_q\cap
 \bigcup_{\substack{3\le n<B\\ n\text{ odd}}}
 \{f(a):a\in A[n](\overline{\mathbf F}_2)\setminus\{0\},
                         \ f(a)\text{ finite and defined}\}\right).
                                                               \tag{AG.33.6}
\]
Here $f(a)$ means the value of the pencil morphism
$A\setminus\{0\}\to\mathbf P^1$, including extension across
removable zeros of the displayed numerator and denominator.
This is an exact algebraic definition using the displayed
rational function on the fixed elliptic surface $E\times E$.
It specifies a finite set, without an assertion that the
large set has been enumerated. The incidence bound gives
\[
 \#\mathcal B\le4+
       \sum_{\substack{3\le n<B\\ n\text{ odd}}}(n^4-1)<4+B^5<q.
\]
Every $t\in\mathbf F_q\setminus\mathcal B$ defines a
smooth curve (AG.33.1), and no odd conductor allowed by
AG.32 can supply a primitive-zero cover. Larger conductors
are impossible by that same bound. This proves the
asserted counterexample existence over the stated finite
field, with no restriction on the field of definition of
the hypothetical cover.

The exact counting certificate actually gives
\[
 \#(\mathbf F_q\setminus\mathcal B)\ \ge\
       340282354541538070615576657943687331836>0.
\]
The standard-library program
\path{code/supersingular_primitive_counterexample.py}
verifies the polynomial identities of (AG.33.2)--(AG.33.4),
the Artin--Schreier coefficient expansion, the elementary
integer bounds, and the exact odd fourth-power sum.
Its complete output is retained in the companion JSON.
The program is a finite algebra and nonemptiness certificate,
not an enumeration of an individual parameter or a test
of infinitely many covers.

Finally, $J_{C_t}$ is isogenous to $E^2$ and hence has
geometric two-rank zero. R1 therefore already excludes a
collision on each $C_t$. The two supersingular elliptic
maps (AG.33.2) jointly separate geometric points as well:
$v_t$ has degree two and a connected kernel, since
$J_{C_t}[2](\overline k)=0$. It is therefore radicial on
geometric points, while the Abel map is an
embedding. This is consistent with AG.30, whose hypothesis
requires a nonzero geometric kernel point.
\end{proof}
